%%%%%%%%%%%%%%%%%%%%%%%%%%%%%%%%%%%%%%%%%%%%%%%%%%%%%%%%%%%%%%%%%%%%%%%%%%%%%%%%%%%%%%%%%%%
% Original author of template for CV (Plasmati Graduate CV v1.0 24/03/2013):
% Alessandro Plasmati (alessandro.plasmati@gmail.com)
% Changes from the template onwards by (latest 30/01/2019):
% Pol del Aguila Pla (poldap@kth.se)
% Downloaded from:
% http://www.LaTeXTemplates.com
% License:
% CC BY-NC-SA 3.0 (http://creativecommons.org/licenses/by-nc-sa/3.0/)
% Important note (PdaP):
% This template needs to be compiled with XeLaTeX.
% The main document font is called Fontin and can be installed in sudo-enabled
% Linux systems by running fixFONTproblem.sh (will internally call sudo when necessary).
%%%%%%%%%%%%%%%%%%%%%%%%%%%%%%%%%%%%%%%%%%%%%%%%%%%%%%%%%%%%%%%%%%%%%%%%%%%%%%%%%%%%%%%%%%

% ----------------------------------------------------------------------------------------
%	PACKAGES AND OTHER DOCUMENT CONFIGURATIONS
% ----------------------------------------------------------------------------------------

% Font and paper size
\documentclass[a4paper,10pt]{article}

\usepackage{setspace}
% Font loading
% \usepackage{fontspec}
% \defaultfontfeatures{Mapping=tex-text}
% Set main font for document
% \setmainfont[SmallCapsFont = Fontin SmallCaps]{Fontin}

% Formatting
% \usepackage{xunicode,xltxtra,url,parskip}

% Coloring
\usepackage[usenames,dvipsnames]{xcolor}

% Margin specification
\usepackage{fullpage}

% Links and other clickable references
\usepackage{hyperref}
% Link colors
\definecolor{linkcolour}{rgb}{0,0.2,0.6}
\hypersetup{colorlinks,breaklinks,urlcolor=linkcolour,linkcolor=linkcolour}

% Costumize section command
\usepackage{titlesec} % Used to customize the \section command
% Text formatting
\titleformat{\section}{\Large\scshape\raggedright}{}{0em}{}[\titlerule]
% Spacing
\titlespacing{\section}{0pt}{3pt}{3pt}

% Insert images
\usepackage{graphicx}

% Footnotes in tables
\usepackage{footnote}
% input of sepcial characters
\usepackage[utf8]{inputenc}

% Tables that can span several pages
\usepackage{longtable}

% Several biblographies[sectionbib/subsectionbib]
\usepackage[subsectionbib]{bibunits}
% chinese
\usepackage[UTF8]{ctex}
% Trademark symbol
\usepackage{textcomp}

% Put references sections in the TOC (PDF and HTML links)
\usepackage[nottoc,numbib]{tocbibind}

\def\myname{刘栋}


\begin{document}

% Remove page numbers
\pagestyle{empty}

% ----------------------------------------------------------------------------------------
%	NAME AND CONTACT INFORMATION
% ----------------------------------------------------------------------------------------

\begin{center}
  \begin{tabular}{lcr}
    \par{\centering{\Huge 刘栋, Ph.D.}\bigskip\par} & & %\includegraphics[width=0.3\textwidth]{../main/pol.jpg} \\
  \end{tabular}
\end{center}

\section{个人信息}

\begin{tabular}{rl}
  \textsc{地址:} & Malvinas Vag 10, Stockholm, Sweden \\
  \textsc{电话:} & +460767103249 \\
  \textsc{邮件:} & \href{mailto:liu-dong@live.cn}{\nolinkurl{liu-dong@live.cn}}, \href{mailto:doli@kth.se}{\nolinkurl{doli@kth.se}} \\
  % \textsc{Google Scholar:} & \href{https://scholar.google.com/citations?user=eK9LoQMAAAAJ&hl=en}{\aiGoogleScholar}\\
  \textsc{个人页面:} & \href{https://people.kth.se/~doli/}{https://people.kth.se/{\textasciitilde}doli/} 或Git IO \href{https://firsthandscientist.github.io}{https://firsthandscientist.github.io}\\
\end{tabular}

% ----------------------------------------------------------------------------------------
% EDUCATION
% ----------------------------------------------------------------------------------------

\section{教育经历}

\begin{tabular}{r|p{12cm}}

  \textsc{2020.11} 	& 博士\\
  \textsc{2016.12} 	& {信息科学与工程实验室, 智能系统系}, {电气工程与计算机学院} \\
                      & {瑞典皇家理工学院，斯德哥尔摩} \\
 		 & {研究方向：统计模型和算法，贝叶斯推理，贝叶斯网络、马尔可夫随机场，机器学习算法设计} \\
                      & {博士论文: \href{https://github.com/FirstHandScientist/phdthesis}{\textbf{Perspectives on Probabilistic Graphical Models}}}，\href{http://kth.diva-portal.org/smash/record.jsf?pid=diva2%3A1476655&dswid=353}{KTH DIVA}。 Opponent：\href{http://users.ics.aalto.fi/harrila/}{Harri Lähdesmäki}, 阿尔托大学, 芬兰。 答辩委员：\href{https://www.chalmers.se/en/staff/Pages/lennart-svensson.aspx}{Lennart Svensson}，查尔莫斯理工大学, 瑞典; \href{https://folk.idi.ntnu.no/helgel/}{Helge Langseth}，挪威科技大学，挪威; \href{https://www.ucl.ac.uk/~ucakima/}{Ioanna Manolopoulou}, 伦敦大学学院，英国。 \\

  \cline{2-2}
  % ------------------------------------------------

  \textsc{2016.3} 	& 工学硕士， \textbf{Top $\mathbf{5\%}$, 2016年上海市优秀毕业生}\\
\textsc{2013.9}  	& {信息与通信工程系，电气工程与计算机学院} \\
                      & {同济大学，上海}\\
& {学位论文：大规模无线网络：认知网络连通性分析与网络优化} \\
  \cline{2-2}
  % ------------------------------------------------

  \textsc{2013 Jul} 	& 工学学士 \\
  \textsc{2009 Sep} & \textbf{Top $\mathbf{1\%}$, 2013年辽宁省优秀毕业生}\\
  & {信息科学与工程学院} \\
                      & {沈阳工业大学，沈阳} \\
  \multicolumn{1}{c}{}\\

\end{tabular}

% ----------------------------------------------------------------------------------------
%	Work
% ----------------------------------------------------------------------------------------

\section{工作经历}
\begin{tabular}{r|p{12cm}}
  \textsc{至今}  & 研究员, 贝叶斯模型和方法 \\
  \textsc{2020.11}  &{电气工程与计算机学院}  \\
& 瑞典皇家理工学院，斯德哥尔摩\\
  \cline{2-2}
  \textsc{2020.1}  & 联合创始人、算法顾问 \\
	\textsc{2018.1}& 人工智能系统方案设计、智能设备供应 \\
		& \href{http://www.siajor.com/}{视距智能SIAJOR Co. Ltd} \\
  \cline{2-2}
    \textsc{2016.12}  & 研究工程师, 无线接入技术  \\
  \textsc{2016.4}  & 华为上海研究所 \\
\end{tabular}


\section{暑期学校}
\begin{tabular}{r|p{13cm}}

  \textsc{2012 Sep}  & 中国科学院自动化研究所, 沈阳，中国。（奖学金由中科院提供） \\
                     
  \multicolumn{1}{c}{}\\
  \textsc{2012 Aug}  & 高知工科大学, 高知, 日本。 （奖学金由高知提供) \\
\end{tabular}

% ----------------------------------------------------------------------------------------
% COMPUTER SKILLS
% ----------------------------------------------------------------------------------------

\section{工作技能(\href{https://github.com/FirstHandScientist}{\textit{GitHub Profile}})}

\begin{tabular}{r|p{12cm}}
	程序语言与平台 & \textsc{Python, Pytorch, Bash, C++，Matlab, \LaTeX}  \vspace{5pt}\\
	开发能力 & 熟悉Linux开发与管理环境，2017-2020期间担任瑞典皇家理工学院ISE计算服务器管理员.
         
\end{tabular}

\section{科研工作介绍}
\begin{tabular}{r|p{13cm}}
\textsc{研究兴趣} & 概率图模型设计和应用 \\
&通用贝叶斯推理/学习算法设计，置信度传播（belief propagation/message passing) \\
& 时序信号模型建模与算法设计 \\
& 随机几何、随机网络、复杂网络模型设计与应用 \\
& 以上模型与算法的应用（语音、图像、性能分析、通信）
\end{tabular}

\subsection{近期项目经历}
\begin{tabular}{r|p{13cm}}
\textsc{2020} & 基于神经网络的概率图模型设计和应用\\
\textsc{2019}& 项目参与单位：瑞典皇家理工学院，卡罗林斯卡学院 \\
              & 承担项目内容：基于神经网络的概率可追溯模型设计和应用，该项目整合最大期望算法、隐马尔可夫模型和基于flow的神经网络模型提出了新型模型和算法; 项目资源平台： \href{https://www.snic.se/}{Swedish National Infrastructure for Computing}（SNIC 2019/3-232)\\
	& 项目完成应用与展示： \href{https://github.com/FirstHandScientist/genhmm}{语音识别}, \href{https://arxiv.org/pdf/1910.13904.pdf}{早产儿败血症检测}, \href{https://github.com/FirstHandScientist/classification-EM-GM}{最大似然图像生成与分类判决}\\
  \cline{2-2}

  \textsc{2019}  & Advanced Message Passing设计和应用\\
  \textsc{2018}& 项目参与单位：瑞典皇家理工学院，华为瑞典研究所。\\
              & 承担项目内容：低复杂度、高效率approximate belief propagation/message passing算法部署、算法设计、以及在MIMO中的应用。项目经费100万克朗+。\\
              & 两名硕士研究生（Scotti, Torres）毕业论文依托于该项目,\textbf{所辅导学生论文获得ICML Workshop GRL＋创新应用提名};部分项目开源于\href{https://github.com/FirstHandScientist/amp}{github repository} \\
  \cline{2-2}
  & 其他近期开源项目见\href{https://people.kth.se/~doli/#/research}{research\&development}\\
  
\end{tabular}


% ----------------------------------------------------------------------------------------
%	SCHOLARSHIPS AND ADDITIONAL INFO
% ----------------------------------------------------------------------------------------
\subsection{其他学者经费\&学生奖学金}
\begin{savenotes}
  \begin{longtable}{r|p{13cm}}
    % General Travel grant: 21500 SEK
    % Ericsson research grant: 20 000 SEK
    % Wallenberg jubilee 10100 SEK
    % SPS 500 USD = 4846 SEK
    % Karl Engvers foundation = 27 000 SEK
    % Current sum: 83446 SEK
    \textsc{2019.11} & 经费$\approx 83k$ 瑞典克朗\\
    \textsc{2017.1} & Karl Engvers 基金会,瑞典, 2020 \\
                      & Knut and Alice Wallenberg 基金会"Jubilee appropriation", 瑞典, 2019 \\
                      & 爱立信科研基金, 瑞典, 2019 \\
                      & Grant from General Travel Foundation, KTH, 瑞典, 2019 \\
                      & Gran from IEEE Signal Processing Society Travel Grant, 2019 \\
    \cline{2-2}
    % National graduate scholarship: 8k * 2.5 = 20k RMB
    % National Scholarship for Postgraduates in 2015: 20k RMB
    % period amount: 40k RMB = 55k SEK
    \textsc{2016 Dec} & 研究生奖学金 $\approx 42k$人民币 \\
    \textsc{2013 Aug} &2013-2016中国研究生国家学业奖学金  \\
                      & 2015年中国研究生国家奖学金\\
    \cline{2-2}
    % National schorlarship: 8k *2 RMB
    % Liaoning provincial: 8k RMB
    % mayor: 6k RMB
    %
    % Special Scholarship from Shenyang University of Technology in 2012 and 2011, respectively: 2K *2 RMB
    % First-class Scholarship from Shenyang University of Technology in 2010: 1.5 K RMB
    % period amount: 35.5K RMB = 48.9K RMB
    \textsc{2013 Jun}	& 学生奖学金 $\approx 35k$ 人民币.\\
    \textsc{2009 Sep}	& 2011年、2012年中国本科生国家奖学金\\
                      & 2012年中国仪器仪表学会一等奖学金\\
                      & 2011沈阳市市长奖学金\\
                      & 2010辽宁省政府奖学金\\
                      & 2010、2011、2012沈阳工业大学奖学金\\
    \multicolumn{2}{c}{} \\

  \end{longtable}
\end{savenotes}

% ----------------------------------------------------------------------------------------
%	PUBLICATIONS
% ----------------------------------------------------------------------------------------
\section{科研论文}

\begin{bibunit}[IEEEtran_Pol]
  \renewcommand\refname{会议}
  \nocite{
    zuxing2020cdc,
    anubhad2020,
    andrea2020,
    liu2020powering,
    honore2020hidden,
    liu2020neural,
    liu2019entropy,
    liu2019alpha,
    liu2019dominant,
    zhang2015effective,
  liu2015node,
  wang2015structural,
  liuenergy,
  ren2015applying,
  ren2015spectram,
  ren2015spectrum1}
  \footnotesize{\putbib[../bibfile]}
\end{bibunit}
\begin{bibunit}[IEEEtran_Pol]
  \renewcommand\refname{期刊}
  \nocite{
    liu2019discontinuous,
    liu2018will,
    liu2015bounds,
    liu2014secondary,
    wang2015rsh}
  \footnotesize{\putbib[../bibfile]}
\end{bibunit}

\begin{bibunit}[IEEEtran_Pol]
  \renewcommand\refname{Preprint}
  \nocite{
    ghosh2020rebust,
    liu2020region,
    liu2020alpha,
    jing2020unfair,
    chatterjee2019ssfn}
  \footnotesize{\putbib[../bibfile]}
\end{bibunit}

\section{发明专利}
 X. Hu, \textbf{D. Liu}, T. Deng, Wireless communication method and apparatus, 覆盖国家：中国（CN108377536A），美国（US20190261448A1），欧盟（EP3528553A1），巴西（BR112019008777A2）。
  




\section{教学经历}
\begin{tabular}{r|p{12cm}}
& 授课经历 \\
		\cline{2-2}
    {\textsc{2020}}& 教学助理，工业界课程(爱立信), {EP232U 深度神经网络} \\
	& 课程内容：时序模型深度网络，深度生成模型（生成对抗网络，变分自动编码器）等\\
  	& 课程提供方：瑞典皇家理工学院.\\
  \cline{2-2}

  \textsc{2020.9}  & 教学助理，研究生课程 \href{https://www.kth.se/student/kurser/kurs/EQ2341?l=en}{EQ2341 模式识别与机器学习} (每学年开课1-2次)\\
  \textsc{2017.9}  & 教学内容：贝叶斯判决与决策，隐马尔可夫模型，变分贝叶斯估计等\\
		& 依托于该课程的研究生项目： 语音识别，手写字体识别，歌曲曲目识别等 \\
& 课程提供方： 瑞典皇家理工学院\\
  \multicolumn{2}{c}{} \\
  % ------------------------------------------------
   {} & 担任硕士研究生论文导师经历 \\
\cline{2-2}
  \textsc{2021} & 任第一导师： Julius Olson, Using machine learning to improve volume estimation and yield prediction in the Swedish forestry industry. 项目合作方 Digital Edge (瑞典)\\
\cline{2-2}
\textsc{2018}  & 任第一导师：Ruichao Qian, Deep-learning-based pose estimation of 3D object, 项目合作方 Gleechi AB (瑞典)\\
\cline{2-2}
  & 任第一导师：任第一导师：Tianyi Li, Exploration of feature extraction algorithm on deepfake detection for videos, 项目合作方 Guizhou YongHong Aviation Machinery（中国）\\
\cline{2-2}
  & 任第一导师：\href{https://www.linkedin.com/in/andrea-scotti-144955184/}{A. Scotti} (topic: graph neural network \& approximate message passing), 现为\href{https://nnaisense.com/}{NNAISENSE} 人工智能工程师（瑞士） \\
 & 指导学位论文：\href{http://kth.diva-portal.org/smash/record.jsf?pid=diva2%3A1479310&dswid=-8558}{Graph Neural Networks and Learned Approximate Message Passing Algorithms for Massive MIMO Detection}， 项目合作方华为瑞典研究所\\
\cline{2-2}
                      &任第二导师： \href{https://www.linkedin.com/in/anubhabghosh/}{A. Ghosh} (topic: neural network \& hidden Markov model), 现为KTH博士生（瑞典） \\
& 指导学位论文：\href{http://www.diva-portal.org/smash/record.jsf?pid=diva2%3A1504139&dswid=6812}{Normalizing Flow based Hidden Markov Models for Phone Recognition} \\
\cline{2-2}
& 任第一导师： \href{https://www.linkedin.cn/in/joao-fernandes-torres/?trk=public_profile_browsemap_profile-result-card_result-card_full-click}{J. Torres} (topic: approximate inference \& message passing), 现为Institute for Systems and Robotics研究员（葡萄牙）\\
& 指导学位论文：\href{http://kth.diva-portal.org/smash/record.jsf?pid=diva2%3A1421287&dswid=7763}{Approximate Inference Low-Complexity MIMO Detection}， 项目合作方华为瑞典研究所\\
\cline{2-2}
& 任第一导师：F. Goncalves Contente \\
& 指导学位论文： \href{http://kth.diva-portal.org/smash/record.jsf?pid=diva2%3A1498763&dswid=-1280}{Hierarchical Clustering based Dynamic Subarrays for Hybrid Beamforming Massive MU-MIMO}，项目合作方华为瑞典研究所

\end{tabular}


% ----------------------------------------------------------------------------------------
% LANGUAGES
% ----------------------------------------------------------------------------------------

\vspace{-10pt}

\section{语言}

\begin{tabular}{r|p{10cm}}

  \textsc{母语} & 中文 \\

  \textsc{工作语言} & English \\

\end{tabular}



% ----------------------------------------------------------------------------------------
%	Awards and Hornors
% ----------------------------------------------------------------------------------------
\section{竞赛经历}

\begin{longtable}[H]{r|p{13.5cm}}
  \hspace{20pt} \textsc{2019} & 第五届中国“互联网+”大学生创新创业大赛国际赛道铜奖 \\
  
  \textsc{2014} & 中国全国研究生数学建模竞赛一等奖，Top \textbf{$\mathbf{2.45\%}$(共4900支队伍)}\\
  
\textsc{2012} & 沈阳市“福田杯”机器人设计竞赛二等奖(区域内寻物机器人的设计与制作)\\
  \textsc{2011} & 全国本科生电子设计竞赛辽宁赛区特等奖(简易的数字信号发生器、分析仪的设计与制作)\\
  \textsc{2011} & 全国本科生数学竞赛辽宁赛区一等奖 \\
\end{longtable}

\section{荣誉奖项}
\begin{longtable}[H]{r|p{13.5cm}}
  \textsc{2016.7} & 2016年上海市优秀毕业生. \\
  \textsc{-}& 2013年辽宁省优秀毕业生. \\
  \textsc{2011.9} & 2012 年辽宁省年度大学生人物提名奖 \\
                    & 2012 年辽宁省三好学生\\
                    & 2011 年辽宁省励志成才优秀大学生 \\
                    %& The Youth Medal of Shenyang University of Technology in 2012 (4 award winners are selected among over 16,000 undergraduates every two years). \\
                    %& The {Second Prize for Outstanding Contribution} to the university in 2011, the {Excellent Student Leader} in 2011, the {Pivot of Merit Students} in 2010, the {Top Ten Students} of school in 2012, 2011 and 2010, respectively, Shenyang University of Technology.
\end{longtable}

% ----------------------------------------------------------------------------------------
%	Contact with the scientific community
% ----------------------------------------------------------------------------------------

% \section{Participation in the scientific community}

% \begin{longtable}[H]{r|p{13.5cm}}
%   \emph{Current}	& Reviewer for \textbf{Elsevier Signal Processing} \\
%   \textsc{2019 Mar} &  \\
%   \multicolumn{2}{c}{} \\

%   \emph{Current}	& Reviewer for the \textbf{IEEE Transactions on Signal Processing} \\
%   \textsc{2015 Aug} &  \\
%   \multicolumn{2}{c}{} \\


%   \textsc{2019 Apr} &  \textsc{12-13 Apr}, Invited participant in the 2019 IEEE Signal Processing Society (SPS) Long Range Planning Meeting\\
%   & \\
%   & Lecture presentation (\textsc{9 Apr}) within the \emph{2019 IEEE $16^{\mathrm{th}}$ International
%   Symposium on Biomedical Imaging} (\textbf{ISBI 2019}), titled
%   \textbf{SpotNet --- Learned iterations for cell detection in image-based immunoassays}, access at \href{https://embs.papercept.net/conferences/scripts/myprogram.pl?ConfID=79&Add=208}{\texttt{embs.papercept.net}}.\\
%   & \footnotesize{Hilton Molino Stucky, Venice, Italy.} \\
%   \multicolumn{2}{c}{} \\



%   \textsc{2019 Jan} & \textsc{7 Jan - 7 Feb}, Research visit at \textbf{Professor Jean-Luc Starck}'s group (\textbf{CosmoStat}). \\
%   & \footnotesize{Department of Astrophysics, \textbf{CEA Paris-Saclay}, Paris, France} \\
%   & \\
%   & \textsc{14 and 28 Jan}, presentations at \textbf{Cosmostat} and \textbf{Parietal}, \textbf{NeuroSpin}, \textbf{INRIA}, respectively, titled \emph{Cell detection by functional inverse diffusion and non-negative group sparsity - Biology, physics, math and engineering}, access at \href{http://www.cosmostat.org/events/cosmoclub/cosmosclub-pol}{\texttt{www.cosmostat.org}} and \href{https://team.inria.fr/parietal/slides-of-pol-del-aguila-plas-talk-available-now-online/}{\texttt{team.inria.fr/parietal}}. \\
%   \multicolumn{2}{c}{} \\

%   \textsc{2018 Jul} & Attendance to the \emph{Thirty-fifth International Conference on Machine
%   Learning} (\textbf{ICML 2018}). \\
%   & \footnotesize{Stockholmsm\"{a}ssan, Stockholm, Sweden.} \\
%   \multicolumn{2}{c}{} \\

%   \textsc{2018 Jun} & Poster presentation within the \emph{SIAM Conference on Imaging Science}
%   (\textbf{SIAM-IS 2018}), titled \textbf{Source localization by spatially
%   variant blind deconvolution}, access at \href{https://www.siam-is18.dm.unibo.it/uploads/store/a7a8b242b168225d0be8998fa373f58b.pdf}{\texttt{www.siam-is18.dm.unibo.it}}. \\
%   & \footnotesize{\textbf{University of Bologna}, Bologna, Italy.}\\
%   \multicolumn{2}{c}{} \\

%   \textsc{2018 Apr} & Poster presentation within the \emph{IEEE International Conference
%   on Acoustics, Speech and Signal Processing} (\textbf{ICASSP 2018}), titled
%   \textbf{Convolutional group-sparse coding and source localization}, access at \href{https://sigport.org/documents/convolutional-group-sparse-coding-and-source-localization}{\texttt{sigport.org}}. \\
%   & \footnotesize{Calgary Talus Convention Centre, Calgary, Alberta, Canada.} \\
%   & \\
%   & \textsc{9 - 13 Apr}, Research visit at \textbf{Professor Stephen P. Boyd's group}. \\
%   & \textsc{10 Apr}, Presentation to the group titled
%   \emph{Cell detection by functional inverse diffusion and non-negative group sparsity}.\\
%   & \footnotesize{Information Systems Laboratory, Department of \textbf{Electrical Engineering,
%   Stanford University}, Stanford, California, United States of America} \\
%   & \\
%   & Poster presentation within the \emph{2018 IEEE $15^{\mathrm{th}}$ International
%   Symposium on Biomedical Imaging} (\textbf{ISBI 2018}), titled
%   \textbf{Cell detection on image-based immunoassays}.\\
%   & \footnotesize{Omni Shoreham Hotel, Washington, D.C., United States of America.} \\
%   \multicolumn{2}{c}{} \\

%   \textsc{2017 Nov} & Lecture presentation within the workshop \emph{Generative models,
%   parameter learning, and sparsity} (VMVW02), titled
%   \textbf{Cell detection by functional inverse diffusion and group sparsity},
%   access at \href{https://downloads.sms.cam.ac.uk/2600830/2600858.mp4}{\texttt{downloads.sms.cam.ac.uk}}. \\
%   & \footnotesize{\textbf{Isaac Newton Institute for Mathematical Sciences},
%   \textbf{University of Cambridge}, Cambridge, United Kingdom.
%   Within the programme \emph{Variational methods and effective
%   algorithms for imaging and vision}.}\\

% \end{longtable}

% ----------------------------------------------------------------------------------------
% TEACHING EXPERIENCE
% ----------------------------------------------------------------------------------------
% \section{Teaching experience}

% \begin{bibunit}[IEEEtran_Pol]
%   \begin{tabular}{r|p{13cm}}

%     \emph{Current}	 & Supervision of Master's and Bachelor's thesis (see below) \\
%     \textsc{2017 Feb}  & \footnotesize{Main supervisor of two master's thesis, \cite{Jones2018} and one
%     ongoing.
%     Designer of two bachelor's thesis projects, attracting $5$
%     different groups of two students. Main supervision for $2$ bachelor's
%     thesis \cite{F2A2018,F2B2018}, and co-supervisor for six,
%     \cite{F3A2018,F3B2018,F3C2018} and three ongoing. } \\
%     \multicolumn{2}{c}{} \\

%     \textsc{2018 Dec}     & Teaching assistance in the course \textbf{EQ2300: Digital Signal Processing} \\
%     \textsc{2014 Sep} & \footnotesize{Taught every year \textsc{Nov -- Dec} by Prof. Joakim Jald\'{e}n and
%     one to three assistants. Approximate numbers: $120\,\mathrm{h}$ of guidance of
%     exercise sessions and lectures, $140\,\mathrm{h}$ of
%     class preparation, $45\,\mathrm{h}$ of course and material development,
%     $100\,\mathrm{h}$ of grading of projects and exams, and $20\,\mathrm{h}$ of private tutoring.} \\

%   \end{tabular}


%   \renewcommand\refname{\normalsize{Supervised theses}}
%   \footnotesize{
%   \putbib[../pubs/bibfile_students]}
% \end{bibunit}


% ----------------------------------------------------------------------------------------
% REFEREE
% ----------------------------------------------------------------------------------------
\section{References}
\begin{longtable}[H]{r|p{13.5cm}}
  Assoc. Prof. & 电气工程与计算机学院\\
  Ragnar Thobaben &瑞典皇家理工学院, 斯德哥尔摩 \\
               & 办公地址: Room C:738, Malvinas vag 10, Stockholm, SE-100 44, Sweden \\
               & Email: ragnart@kth.se \\
               & 电话: +4687908452 \\

  \multicolumn{2}{c}{} \\
  Assoc. Prof. & 电气工程与计算机学院\\
  Saikat Chatterjee &瑞典皇家理工学院，斯德哥尔摩 \\
               & 办公地址: Malvinas vag 10, Stockholm, SE-100 44, Sweden \\
               & Email: sach@kth.se \\
               & 电话: +4687908478 \\

  \multicolumn{2}{c}{} \\

  Prof.  & 电气工程与计算机学院 \\
  Mikael Skoglund &瑞典皇家理工学院，斯德哥尔摩 \\
               & 办公地址: Room A:415, Malvinas vag 10, Stockholm, SE-100 44, Sweden \\
               & Email: skoglund@kth.se \\
               & 电话: +4687908430 \\
	
  \multicolumn{2}{c}{} \\
  Assoc. Prof. & 电子与信息工程学院 \\
  王超 & 同济大学, 上海 \\
               & 办公地址: 上海市嘉定区曹安公路4800号\\
               & Email: chaowang@tongji.edu.cn \\
               & 电话: +86 21 69583671\\
                 
\multicolumn{2}{c}{} \\
  Prof. & 电子与信息工程学院 \\
  刘儿兀 & 同济大学, 上海 \\
               & 办公地址: 上海市嘉定区曹安公路4800号\\
               & Email: erwu.liu@ieee.org or erwuliu@tongji.edu.cn \\
               & 电话: +8615800903275

\end{longtable}



\end{document}
%%% Local Variables:
%%% mode: latex
%%% TeX-master: t
%%% End:
